%==============================================================================
%  ANNEXE - Exploitation et engagements de service
%  SLO, sauvegardes, reversibilite : ce qui rend le systeme livrable
%==============================================================================

%==============================================================================
\chapter{Annexe : Exploitation et engagements de service}
\label{chap:exploitation}
%==============================================================================
Un système livré sans engagement de service mesurable n'est pas exploitable :
faute d'objectif, aucune dégradation ne peut être qualifiée d'incident. Cette
annexe fixe les objectifs de service, la politique de sauvegarde et les
conditions de réversibilité. Elle sert de référence aux critères de recette du
chapitre~\ref{chap:delais}.

%------------------------------------------------------------------------------
\section{Objectifs de service (SLO)}
%------------------------------------------------------------------------------
Les objectifs sont dimensionnés pour un déploiement sur \textbf{serveur unique},
conformément à la décision d'exploitation retenue. Ils sont délibérément
réalistes : promettre 99{,}99~\% sans architecture redondante serait un
engagement intenable.

\begin{xltabular}{\textwidth}{>{\bfseries}p{4.2cm} p{3.4cm} X}
\toprule
\textbf{Indicateur} & \textbf{Objectif} & \textbf{Précision} \\
\midrule
\endhead
Disponibilité & 99{,}5~\% mensuel & Environ 3~h~40 d'indisponibilité tolérée par mois, hors maintenances planifiées. \\
\addlinespace
Maintenance planifiée & 4~h/mois maximum & Annoncée 7 jours à l'avance, hors saison de production quand c'est possible. \\
\addlinespace
Latence (p95) & $<$ 500~ms & Sur les écrans de consultation, à charge nominale. \\
\addlinespace
Latence (p99) & $<$ 1~500~ms & Tolère les requêtes analytiques lourdes. \\
\addlinespace
Taux d'erreur & $<$ 1~\% & Réponses 5xx rapportées au total des requêtes. \\
\addlinespace
\textbf{RTO} (reprise) & $<$ 4~h & Délai maximal de remise en service après incident majeur. \\
\addlinespace
\textbf{RPO} (perte de données) & $<$ 1~h & Perte maximale acceptable, garantie par la fréquence de sauvegarde. \\
\bottomrule
\end{xltabular}

\begin{encadre}
\textbf{Charge nominale de référence :} 200 utilisateurs simultanés en pic de
saison, 10~000 ruches instrumentées, environ 1~million de mesures ingérées par
jour. Les objectifs de latence ne valent que dans ces limites ; au-delà, le
dimensionnement doit être révisé.
\end{encadre}

%------------------------------------------------------------------------------
\section{Sauvegarde et restauration}
%------------------------------------------------------------------------------
\subsection{Politique}
\begin{itemize}
  \item \textbf{Sauvegarde continue} des journaux de transaction PostgreSQL
        (\emph{WAL archiving}), garantissant le RPO d'une heure.
  \item \textbf{Sauvegarde complète quotidienne}, chiffrée au repos.
  \item \textbf{Externalisation} sur un support distinct du serveur de
        production --- une sauvegarde stockée sur la machine qu'elle protège ne
        protège de rien.
  \item \textbf{Rétention} : 7 sauvegardes quotidiennes, 4 hebdomadaires,
        12 mensuelles.
  \item Les \textbf{photos de visite} et le fichier \texttt{ConfigZumm.ini}
        entrent dans le périmètre de sauvegarde au même titre que la base.
\end{itemize}

\subsection{Test de restauration}
\begin{encadre}
\textbf{Une sauvegarde jamais restaurée n'est pas une sauvegarde.} La
restauration complète est testée \textbf{mensuellement} sur un environnement
dédié, et le compte rendu est archivé. Un test de restauration réussi est un
critère de recette bloquant.
\end{encadre}

Le test vérifie trois points : l'intégrité des données restaurées, le respect
du RTO annoncé, et la capacité d'un exploitant à mener la procédure en
s'appuyant sur le seul runbook.

%------------------------------------------------------------------------------
\section{Supervision et alertes}
%------------------------------------------------------------------------------
\begin{xltabular}{\textwidth}{>{\bfseries}p{4.2cm} p{3.4cm} X}
\toprule
\textbf{Signal} & \textbf{Seuil d'alerte} & \textbf{Action attendue} \\
\midrule
\endhead
Indisponibilité applicative & 2 échecs consécutifs & Intervention immédiate \\
Espace disque & $>$ 80~\% & Extension ou purge sous 48~h \\
Latence p95 & $>$ 500~ms sur 15~min & Analyse de la requête en cause \\
Taux d'erreur 5xx & $>$ 1~\% sur 5~min & Intervention immédiate \\
Échec de sauvegarde & 1 occurrence & Intervention le jour même \\
Certificat TLS & expiration $<$ 30~j & Renouvellement \\
Retard d'ingestion capteurs & $>$ 1~h & Vérification du courtier MQTT \\
\bottomrule
\end{xltabular}

La supervision s'appuie sur \textbf{OpenTelemetry}, Prometheus et Grafana,
conformément à la pile de l'annexe~\ref{chap:technologies}.

%------------------------------------------------------------------------------
\section{Réversibilité}
%------------------------------------------------------------------------------
Le client doit pouvoir reprendre son système sans dépendre du prestataire.
Sont remis à la livraison, et à toute demande ultérieure :

\begin{itemize}
  \item l'\textbf{export complet des données} dans un format ouvert et
        documenté (SQL et CSV), incluant les photos et les fichiers joints ;
  \item le \textbf{code source} et l'historique de versions ;
  \item la \textbf{documentation d'exploitation} (runbook) : démarrage, arrêt,
        restauration, rotation des certificats, lecture des alertes ;
  \item les \textbf{schémas de base} et les scripts de migration ;
  \item la \textbf{procédure de redéploiement} depuis un dépôt vierge, vérifiée
        lors de la recette.
\end{itemize}

\begin{encadre}
\textbf{Vérification :} la réversibilité se démontre, elle ne se déclare pas.
Un redéploiement à blanc depuis le dépôt, sur une machine neuve, est réalisé
pendant la recette.
\end{encadre}

%------------------------------------------------------------------------------
\section{Transfert de compétences}
%------------------------------------------------------------------------------
Deux journées de transfert sont prévues à la livraison, couvrant
l'administration courante, la lecture de la supervision, la procédure de
restauration et le diagnostic de premier niveau. Le runbook sert de support et
reste la référence après le transfert.
